% from ega1-10-fr.tex
\oldpage[I]{209}

  \item[\rm{(ii)}] Soient $\mathfrak{X}$, $\mathfrak{S}$,
    $\mathfrak{S}'$ trois préschémas formels localement noethériens (resp.
    noethériens), $f:\mathfrak{X}\to\mathfrak{S}$,
    $g:\mathfrak{S}'\to\mathfrak{S}$ deux morphismes. Si $f$ est de type
    fini, $\mathfrak{X}\times_{\mathfrak{S}}\mathfrak{S}'$ est localement
    noethérien (resp. noethérien) et est de type fini sur $\mathfrak{S}'$.
  \item[\rm{(iii)}] Soient $\mathfrak{S}$ un préschéma formel localement
    noethérien, $\mathfrak{X}'$, $\mathfrak{Y}'$ deux
    $\mathfrak{S}$-préschémas formels localement noethériens tels que
    $\mathfrak{X}'\times_{\mathfrak{S}}\mathfrak{Y}'$ soit localement
    noethérien. Si $\mathfrak{X}$, $\mathfrak{Y}$ sont des
    $\mathfrak{S}$-préschémas formels localement noethériens,
    $f:\mathfrak{X}\to\mathfrak{X}'$,
    $g:\mathfrak{Y}\to\mathfrak{Y}'$ deux $\mathfrak{S}$-morphismes de
    type fini, $\mathfrak{X}\times_{\mathfrak{S}}\mathfrak{Y}$ est
    localement noethérien et $f\times_{\mathfrak{S}}g$ est un
    $\mathfrak{S}$-morphisme de type fini.
\end{enumerate}
\end{proposition}

(iii) se déduit de (i) et (ii) par le raisonnement formel de
(\hyperref[I.3.5.1-fr]{3.5.1}) et il suffit donc de prouver (i) et (ii).

Soient $\mathfrak{X}$, $\mathfrak{Y}$, $\mathfrak{Z}$ trois préschémas
formels localement noethériens, $f:\mathfrak{X}\to\mathfrak{Y}$,
$g:\mathfrak{Y}\to\mathfrak{Z}$ deux morphismes de type fini. Si
$\mathscr{L}$ est un Idéal de définition de $\mathfrak{Z}$,
$\mathscr{K}=g^*(\mathscr{L})\mathscr{O}_{\mathfrak{Y}}$ en est un pour
$\mathfrak{Y}$ et
$\mathscr{J}=f^*(g^*(\mathscr{L}))\mathscr{O}_{\mathfrak{X}}$ en est un
pour $\mathfrak{X}$. Posons
$X_0=(\mathfrak{X},\mathscr{O}_{\mathfrak{X}}/\mathscr{J})$,
$Y_0=(\mathfrak{Y},\mathscr{O}_{\mathfrak{Y}}/\mathscr{K})$,
$Z_0=(\mathfrak{Z},\mathscr{O}_{\mathfrak{Z}}/\mathscr{L})$ et soient
$f_0:X_0\to Y_0$, $g_0:Y_0\to Z_0$ les morphismes correspondant à $f$ et
$g$. Comme par hypothèse $f_0$ et $g_0$ sont de type fini, il en est de même
de $g_0\circ f_0$ (\hyperref[I.6.3.4-fr]{6.3.4}) qui correspond à
$g\circ f$~; donc $g\circ f$ est de type fini par
(\hyperref[I.10.13.1-fr]{10.13.1}).

Sous les conditions de (ii), $\mathfrak{S}$ (resp. $\mathfrak{X}$,
$\mathfrak{S}'$) est limite inductive d'une suite $(S_n)$ (resp. $(X_n)$,
$(S'_n)$) de préschémas localement noethériens et on peut supposer
(\hyperref[I.10.13.1-fr]{10.13.1}) que
$X_m=X_n\times_{S_n}S_m$ pour $m\leq n$. Le préschéma formel
$\mathfrak{X}\times_{\mathfrak{S}}\mathfrak{S}'$ est alors limite inductive
des préschémas $X_n\times_{S_n}S'_n$
(\hyperref[I.10.7.4-fr]{10.7.4}), et on a
\[
  X_m\times_{S_m}S'_m
  =(X_n\times_{S_n}S_m)\times_{S_m}S'_m
  =(X_n\times_{S_n}S'_n)\times_{S'_n}S'_m.
\]
En outre, $X_0\times_{S_0}S'_0$ est localement noethérien puisque $X_0$
est de type fini sur $S_0$ (\hyperref[I.6.3.8-fr]{6.3.8}). On en conclut
d'abord (\hyperref[I.10.12.3.1-fr]{10.12.3.1}) que
$\mathfrak{X}\times_{\mathfrak{S}}\mathfrak{S}'$ est localement
noethérien~; en outre, comme $X_0\times_{S_0}S'_0$ est de type fini sur
$S'_0$ (\hyperref[I.6.3.8-fr]{6.3.8}), il résulte de
(\hyperref[I.10.12.3.1-fr]{10.12.3.1}) et de
(\hyperref[I.10.13.1-fr]{10.13.1}) que
$\mathfrak{X}\times_{\mathfrak{S}}\mathfrak{S}'$ est de type fini sur
$\mathfrak{S}'$, ce qui achève de prouver (ii) (l'assertion relative aux
préschémas noethériens étant conséquence immédiate de
(\hyperref[I.6.3.8-fr]{6.3.8})).

\begin{corollary}[10.13.6]
\label{I.10.13.6-fr}
Sous les hypothèses de (\hyperref[I.10.9.9-fr]{10.9.9}), si $f$ est un
morphisme de type fini, il en est de même de son prolongement $\widehat f$
aux complétés.
\end{corollary}

\subsection{Sous-préschémas fermés des préschémas formels.}
\label{subsection:I.10.14-fr}

\begin{proposition}[10.14.1]
\label{I.10.14.1-fr}
Soient $\mathfrak{X}$ un préschéma formel localement noethérien,
$\mathscr{A}$ un faisceau d'idéaux cohérent de
$\mathscr{O}_{\mathfrak{X}}$. Si $\mathfrak{Y}$ est le support (fermé) de
$\mathscr{O}_{\mathfrak{X}}/\mathscr{A}$, l'espace topologiquement annelé
$(\mathfrak{Y},(\mathscr{O}_{\mathfrak{X}}/\mathscr{A})|\mathfrak{Y})$ est
un préschéma formel localement noethérien, qui est noethérien si
$\mathfrak{X}$ l'est.
\end{proposition}

Notons que $\mathscr{O}_{\mathfrak{X}}/\mathscr{A}$ est cohérent en vertu de
(\hyperref[I.10.10.3-fr]{10.10.3}) et
(\hyperref[0.5.3.4-fr]{0, 5.3.4}), donc son support $\mathfrak{Y}$ est fermé
(\hyperref[0.5.2.2-fr]{0, 5.2.2}). Soit $\mathscr{J}$ un Idéal de définition
de $\mathfrak{X}$, et soit
$X_n=(\mathfrak{X},\mathscr{O}_{\mathfrak{X}}/\mathscr{J}^{n+1})$~; le
faisceau d'anneaux $\mathscr{O}_{\mathfrak{X}}/\mathscr{A}$ est limite
projective des faisceaux
$\mathscr{O}_{\mathfrak{X}}/(\mathscr{A}+\mathscr{J}^{n+1})
=(\mathscr{O}_{\mathfrak{X}}/\mathscr{A})
\otimes_{\mathscr{O}_{\mathfrak{X}}}
(\mathscr{O}_{\mathfrak{X}}/\mathscr{J}^{n+1})$
(\hyperref[I.10.11.3-fr]{10.11.3}), qui ont tous pour support
$\mathfrak{Y}$. Le faisceau
$(\mathscr{A}+\mathscr{J}^{n+1})/\mathscr{J}^{n+1}$ est un
$\mathscr{O}_{\mathfrak{X}}$-Module cohérent, puisque
$\mathscr{J}^{n+1}$ est cohérent, donc
$(\mathscr{A}+\mathscr{J}^{n+1})/\mathscr{J}^{n+1}$ est aussi un
$(\mathscr{O}_{\mathfrak{X}}/\mathscr{J}^{n+1})$-Module cohérent
(\hyperref[0.5.3.10-fr]{0, 5.3.10})~; si $Y_n$ est le sous-préschéma fermé
de $X_n$ défini par ce faisceau d'idéaux, il est immédiat que
$(\mathfrak{Y},(\mathscr{O}_{\mathfrak{X}}/\mathscr{A})|\mathfrak{Y})$

